\documentclass{beamer}
\usetheme{Madrid}

% --- Packages and operators
\usepackage{amsmath,amssymb}
\usepackage{mathtools}
\usepackage{bm}

% --- Macros and operators (standardized)
\newcommand{\inv}[1]{#1^{-1}}
\DeclareMathOperator{\GL}{GL}
\DeclareMathOperator{\SL}{SL}
\DeclareMathOperator{\PSL}{PSL}
\DeclareMathOperator{\PGL}{PGL}
\DeclareMathOperator{\Sym}{Sym}
\DeclareMathOperator{\im}{im}
\newcommand{\PGammaL}{\mathrm{P\Gamma L}}
\newcommand{\Oodd}{\mathbf{O}_{2'}} % largest normal odd-order subgroup
\newcommand{\aand}{\quad {\textrm{and}} \quad} % largest normal odd-order subgroup

\title[CFSG]{The Classification of the Finite Simple Groups}
\subtitle{Lecture Series at Tsinghua University}
\author{Damiano Testa}
\institute[Warwick]{Mathematics Institute, University of Warwick}
\date{}

\begin{document}

\frame{\titlepage}

%------------------------------------------------------------
\begin{frame}{Introduction and notation}

Throughout this lecture, we will be mostly talking about \emph{groups}.

\bigskip

A \emph{group} is a set $G$ with a binary operation $(\_ \cdot \_)\colon G\times G\to G$ (often written by juxtaposition) such that:
\begin{itemize}
  \item \textbf{associativity}: for all $a,b,c\in G$, the identity $(ab)c=a(bc)$ holds;
  \item \textbf{identity}: there exists $e\in G$ with $ae=ea=a$ for all $a$;
  \item \textbf{inverses}: for each $a$ there exists $a^{-1}$ with $a^{-1}a=aa^{-1}=e$.
\end{itemize}

\bigskip

An \emph{involution} is an element $i \in G$ with $i \neq e$ and $i^2=e$.

\smallskip

Involutions play a major role in the classification of finite simple groups.

\end{frame}

%------------------------------------------------------------
\begin{frame}{Examples}

\begin{itemize}
  \item
  $(\mathbb{Z}, +)$ and $(\mathbb{Z} / n \mathbb{Z}, +)$

  In both cases, the identity is $0$ and the inverse of $a$ is $-a$.

  \medskip

  These groups are \emph{abelian}: $a+b=b+a$.
  \medskip
  \item
  Let $A$ be any set (e.g., $A=\{1,2,\dots,n\}$).

  \medskip

  The group $\Sym(A)$ is the group of bijections of $A$ with itself.

  \medskip
  \item
  If $R$ is a ring (associative, not necessarily commutative), then the units of $R$ form a group under multiplication.

  \medskip

  E.g., $\GL_n(k)$ is the group of invertible $n \times n$ matrices over a field~$k$.
\end{itemize}

\end{frame}

%------------------------------------------------------------
\begin{frame}{Conjugacy and normal subgroups}
Let $G$ be a group and $g\in G$.
The \emph{conjugate} of $x\in G$ by $g$ is $g x \inv{g}$.

\bigskip

The conjugate of a \emph{subgroup} $H$ of $G$ is the subgroup
\[
gHg^{-1}= \{gh\inv{g} \mid h \in H\}.
\]

\medskip

A subgroup $N$ of $G$ is \emph{normal} if it is stable under conjugation by all of $G$:
\[
gNg^{-1}= \{gn\inv{g} \mid n \in N\}=N \quad \text{for all } g\in G.
\]

\medskip

If $N$ is normal in $G$, then there is a quotient group $G/N$ and a quotient map $q \colon G \to G/N$ mapping products to products.

\medskip

The map $q$ is a surjective group homomorphism and the only elements of $G$ mapping to the identity are the elements of $N$.

\bigskip

Informally, $G$ is ``built from'' $N$ and $G/N$.

\end{frame}

%------------------------------------------------------------
\begin{frame}{$p$-groups and Sylow subgroups}

A finite group $P$ is a \emph{$p$-group} if $|P|$ is a power of a prime $p$.

\bigskip

Let $G$ be a finite group and let $|G|=p^n m$ with $p\nmid m$.

\bigskip

A subgroup $P\le G$ of order $p^n$ is called a \emph{Sylow $p$-subgroup} of~$G$.

\bigskip

\begin{itemize}
  \item Sylow $p$-subgroups always exist.
  \item Every $p$-subgroup of $G$ is contained in a Sylow $p$-subgroup.
  \item Any two Sylow $p$-subgroups are conjugate in $G$.
\end{itemize}

\bigskip

For instance, we will see Sylow 2-subgroups in a later slide.

\end{frame}

%------------------------------------------------------------
\begin{frame}{Solvable groups and simple groups}

\begin{block}{Definition (Solvable group)}
A group $G$ is \emph{solvable} if there is a finite chain
\[
  G_0 = \{e\} \le G_1 \le \cdots \le G_n = G
\]
of subgroups, where, for each $i \in \{0, \ldots , n - 1\}$, the group $G_i$ is normal in $G_{i+1}$ and the quotient $G_{i+1}/G_i$ is abelian.
\end{block}

{\emph{Solvable}} groups are groups whose (recursive) building blocks are abelian.

\bigskip

In some sense, the opposite of solvable is simple.

\bigskip

\begin{block}{Definition (Simple group)}
  A nontrivial group is {\emph{simple}} if its only normal subgroups are $\{e\}$ and itself.
\end{block}

\end{frame}

%------------------------------------------------------------
\begin{frame}{The statement of the classification}

What are all the finite simple groups, then?

\bigskip

\begin{block}{The classification of finite simple groups}
The finite simple groups are:

\medskip

\begin{itemize}
  \item the cyclic groups of prime order; \hfill (related to $\mathbb{Z}/n\mathbb{Z}$)
  \item the alternating groups; \hfill (related to $\Sym \{1, \ldots , n\}$)
  \item the groups of Lie type; \hfill (related to $\GL_n(k)$)
  \item the 26 sporadic groups. % \hfill (related to $\mathbb{Z}/n\mathbb{Z}$).
\end{itemize}
\end{block}

In the proof, these are the \emph{known groups}.


\bigskip

The goal is to show that no other finite simple groups exist.
\end{frame}

%------------------------------------------------------------
\begin{frame}{Minimal counterexample}
Let $G$ a finite simple group of minimum order among the ones that have not yet been classified.

\bigskip

By minimality, every proper subgroup of $G$ has all composition factors among the known simple groups.

\bigskip

The task is to deduce that $G$ itself is known.

\end{frame}

%------------------------------------------------------------
\begin{frame}{The Feit--Thompson Theorem}
A key starting point is the Feit--Thompson (Odd Order) Theorem.

\begin{block}{Theorem (Feit--Thompson, 1963)}
Every finite group of odd order is solvable.
\end{block}

\bigskip

We may assume our minimal unknown finite simple group $G$ is not cyclic of prime order.

\bigskip

Hence $G$ is not solvable.

\bigskip

The Feit--Thompson Theorem implies that $|G|$ is even.

\bigskip

In particular, $G$ contains elements of order $2$.
\end{frame}

%------------------------------------------------------------
\begin{frame}{Centralisers of involutions}
The Feit--Thompson Theorem guarantees that any non-cyclic finite simple group $G$ contains an involution.

\bigskip

This is the starting point of most classification results.

\bigskip

Many steps analyse the structure of \emph{centralisers} of involutions in $G$.

\bigskip

The \emph{centraliser} of an element $x \in G$ is the subgroup
\[
C_G(x)=\{g\in G \mid gx=xg\}.
\]

\bigskip

Centralisers of elements of prime order also receive a lot of attention.

\bigskip

Information about these centralisers often forces $G$ into a short list of possibilities.

\end{frame}

%------------------------------------------------------------
\begin{frame}{Gorenstein--Walter 1964}
\begin{block}{Theorem (Gorenstein--Walter, 1964)}
A finite simple group whose Sylow $2$--subgroups are dihedral is isomorphic to one of:
\begin{itemize}
  \item $PSL_2(q)$ with $q$ odd, $q>3$;
  \item the alternating group $A_7$.
\end{itemize}
\end{block}

\bigskip

This is a consequence of a more general characterisation.
\end{frame}

%------------------------------------------------------------
\begin{frame}{Gorenstein--Walter 1964 (general form)}

For a finite group $G$, denote by $\Oodd(G)$ the maximal normal subgroup of $G$ of odd order.

\bigskip

\begin{block}{Theorem (Gorenstein--Walter, 1964)}
Let $G$ be a finite group whose Sylow $2$-subgroups are dihedral.

\medskip

Then $G/\Oodd(G)$ is isomorphic to:

\smallskip

\begin{itemize}
  \item a subgroup of $\PGammaL_2(q)$ containing $\PSL_2(q)$, for some odd $q$;
  \item the alternating group $A_7$;
  \item a $2$-group.
\end{itemize}
\end{block}

In particular, the simple groups with dihedral Sylow $2$-subgroups are precisely $\PSL_2(q)$ (odd $q>3$) and $A_7$.
\end{frame}

%------------------------------------------------------------
\begin{frame}{Observations}
\begin{enumerate}
  \item
  The theorem classifies \emph{all} finite groups with dihedral Sylow $2$-subgroups, not just simple ones.

  This is crucial for inductive and reductive steps.
  \item
  The simple groups that occur are exactly $\PSL_2(q)$ (odd $q>3$) and $A_7$

  No other simple groups have dihedral Sylow $2$-subgroups.
  \item
  Groups ``close to simple'' are important.

  Here, subgroups of $P\Gamma L_2(q)$ containing $PSL_2(q)$ play a key role.
  \item
  The subgroup $\Oodd(G)$ (the largest normal subgroup of odd order) is not constrained by the conclusion; it is solvable (by Feit--Thompson).
\end{enumerate}

\bigskip

A later result of Alperin--Brauer--Gorenstein treats the cases where the Sylow $2$-subgroups are quasi-dihedral or wreathed.

\end{frame}


%------------------------------------------------------------
\begin{frame}{The known groups (refined)}

It is often helpful to separate Lie-type groups by the characteristic of the defining field:
\bigskip

\begin{itemize}
  \item cyclic groups of prime order;
  \item alternating groups;
  \item groups of Lie type over fields of \emph{odd characteristic};
  \item groups of Lie type over fields of \emph{characteristic $2$};
  \item the 26 sporadic groups.
\end{itemize}

\bigskip

Characteristic $2$ is separated because involutions behave differently.
\end{frame}

%------------------------------------------------------------
\begin{frame}{Centralisers of involutions in $\SL_3(q)$}

\centering
\begin{tabular}{c|c}
\parbox{0.45\textwidth}{\centering \textbf{Odd characteristic ($q$ odd)}} &
\parbox{0.45\textwidth}{\centering \textbf{Characteristic $2$ ($q=2^f$)}} \\[5pt]
\hline

\parbox{0.45\textwidth}{\centering Involutions are \textbf{semisimple}} &
\parbox{0.45\textwidth}{\centering Involutions are \textbf{unipotent}} \\[5pt]

\parbox{0.45\textwidth}{
\[
t=\begin{pmatrix}
-1 & 0 & 0\\
0 & -1 & 0\\
0 & 0 & 1
\end{pmatrix}
\]} &
\parbox{0.45\textwidth}{
\[
u=\begin{pmatrix}
1 & 1 & 0\\
0 & 1 & 0\\
0 & 0 & 1
\end{pmatrix}
\]} \\[25pt]
%\hline

\parbox{0.45\textwidth}{\centering \textbf{Centraliser}} &
\parbox{0.45\textwidth}{\centering \textbf{Centraliser}}\\[-10pt]
\parbox{0.45\textwidth}{
\[
C_{\mathrm{SL}_3(q)}(t)\subset
\begin{pmatrix}
GL_2(q) & 0\\
0 & *
\end{pmatrix}
\]} &
\parbox{0.45\textwidth}{
\[
C_{\mathrm{SL}_3(q)}(u)\subset
\begin{pmatrix}
* & * & *\\
0 & * & *\\
0 & 0 & *
\end{pmatrix}
\]} \\
%\hline

\parbox{0.45\textwidth}{\centering \textbf{Reductive (Lie type)}} &
\parbox{0.45\textwidth}{\centering \textbf{Solvable (Borel type)}} \\
\end{tabular}

\end{frame}

%------------------------------------------------------------
\begin{frame}{Families of groups of Lie type}
Of course, the ``groups of Lie type'' are also not an indivisible block.

\medskip

They naturally subdivide into subclasses.

\bigskip

\begin{itemize}
  \item \textbf{Classical:} linear, symplectic, even/odd orthogonal
  \item \textbf{Exceptional (untwisted):} $G_2$, $F_4$, $E_6$, $E_7$, $E_8$
  \item \textbf{Twisted:} $^{2}A_n$ (unitary), $^{2}D_n$ (including triality $^{3}D_4$), $^{2}E_6$, and the Suzuki $^{2}B_2(2^{2m+1})$, Ree $^{2}G_2(3^{2m+1})$, $^{2}F_4(2^{2m+1})$ families.
\end{itemize}

\bigskip

In classification arguments, each one must either occur or be ruled out.
\end{frame}

%------------------------------------------------------------
\begin{frame}{The 26 sporadic groups}
Each sporadic group comes with distinctive local and global features:

\bigskip

\begin{itemize}
  \item Sylow subgroups,
  \item subgroup lattices,
  \item centralisers of elements,
  \item geometries,
\end{itemize}
and so on.

\bigskip

Every sporadic group arises somewhere in the classification and must be handled individually.
\end{frame}

%------------------------------------------------------------
\begin{frame}{Case analysis and discovery}
The proof consists of a massive case analysis.

\bigskip

Each case proves a classification theorem under additional hypotheses.

\bigskip

Once all options are covered, the full list of finite simple groups emerges.

\bigskip

Several sporadic groups were discovered through such local analyses.

\bigskip

These relations reveal that the classification is a tightly woven network of interconnected structures, and not a list of isolated objects.

\bigskip

These internal relations are instrumental in the proof.

\end{frame}

%------------------------------------------------------------
\begin{frame}{Computer calculations}
Besides pen-and-paper arguments, many parts of the classification require substantial computer calculations verifying properties of finite simple groups.

\bigskip

Many of these computations have not been thoroughly verified.

\bigskip

Formally checking these computations is a challenge of a different nature than formalising proofs intended for humans.
\end{frame}

%------------------------------------------------------------
\begin{frame}{Humans, AI, and computation}
Many parts of the classification are driven by broad organizing principles.

\bigskip

These ideas shape the whole proof and provide insight and direction.

\bigskip

Yet the endgame often reduces to case-specific arguments with limited transferability.

\bigskip

This mixture can possibly be leveraged to parallelise automated formalisation.

\bigskip

Even so, the formalisation of the proof of the classification of finite simple groups remains a daunting mathematical, software engineering and technological challenge.
\end{frame}

%------------------------------------------------------------
\begin{frame}{Conclusion}
The CFSG is probably the longest proofs in mathematics.

\bigskip

Formalizing it raises multiple challenges:
\begin{itemize}
  \item the statement alone already involves complex mathematical objects (families of groups of Lie type, sporadic groups);
  \item assembling extensive background results across group theory, representation and character theory, algebraic groups, Lie theory, finite geometry, homological algebra, etc.;
  \item managing substantial computational input alongside human-readable arguments,
\end{itemize}
to name a few.

\bigskip

This does not contemplate the possibility that a gap is found and a fix is needed!
\end{frame}

%------------------------------------------------------------
\begin{frame}
\begin{center}
  {\huge Thank you!}\\[30pt]
  {\huge Questions?}
\end{center}

\end{frame}

%------------------------------------------------------------
\end{document}
